% section_iv_corrected.tex
% Corrected Section IV — Harmonic Expansion and Time Rigidity
%
% Fixes applied (referee report):
%   F4.3  SL₂(𝔽ₚ) is finite → amenable: replaced by expander-family framing
%   F4.5  Topological entropy of finite det. system = 0: entropy dropped; spectral gap used
%   F4.4  NCTO circularity + undefined "continuous time": replaced by Gromov–Hausdorff statement
%   F4.6  Factor 24 in M_p = 24p unexplained: explained via CRT splitting
%   F4.2  Generator μ ∉ SL₂(ℤ) (det = −1): replaced by Weyl element w ∈ SL₂(ℤ)
%   F4.1  "Manifold" for finite sets: terminology replaced throughout
%
% Mathematical framework:
%   Tits alternative (1972): G ≤ GL₂(ℚ) finitely generated → virtually solvable or contains F₂
%   Strong approximation (Matthews–Vaserstein–Weisfeiler): Zariski-dense G → ρ_p surjective
%   Bourgain–Gamburd (2008): Zariski-dense lift → expander family with uniform spectral gap
%   Gromov–Hausdorff: expander family with uniform gap ε → no 1D smooth limit

% ─── Macros (assume these are defined in main preamble) ──────────────────────
% \newcommand{\GG}{\mathcal{G}}
% \newcommand{\HH}{\mathcal{H}}
% \newcommand{\VV}{\mathcal{V}}
% \newcommand{\FF}{\mathbb{F}}
% \newcommand{\ZZ}{\mathbb{Z}}
% \newcommand{\QQ}{\mathbb{Q}}
% \newcommand{\RR}{\mathbb{R}}
% \newcommand{\NN}{\mathbb{N}}

\section{Harmonic Expansion and Time Rigidity}
\label{sec:harmonic}

\subsection{Generator-Relative Setup}
\label{sec:harmonic:setup}

Let $\Gamma \subset \mathrm{SL}_2(\mathbb{Z})$ be a finite symmetric set of generators
($\Gamma = \Gamma^{-1}$, $|\Gamma| = d$),
and let $G = \langle \Gamma \rangle \leq \mathrm{SL}_2(\mathbb{Z})$ be the group they generate.

For a prime $p > 3$, let $\rho_p \colon \mathrm{SL}_2(\mathbb{Z}) \to \mathrm{SL}_2(\mathbb{F}_p)$
be the reduction-mod-$p$ homomorphism, and define the Cayley graph
\[
  \mathcal{G}_p \;=\; \mathrm{Cay}\!\bigl(\rho_p(G),\;\rho_p(\Gamma)\bigr).
\]

\begin{definition}[Reachability time and spectral gap]
  \label{def:time_gap}
  The \emph{reachability time} from $x$ to $y$ in $\mathcal{G}_p$ is
  \[
    T_\Gamma(x \to y) \;=\; \min\bigl\{ k : y = g_k \cdots g_1 x,\; g_i \in \rho_p(\Gamma) \bigr\}.
  \]
  The \emph{diameter} is $\mathrm{diam}(\mathcal{G}_p) = \max_{x,y} T_\Gamma(x \to y)$.
  The \emph{spectral gap} $\gamma_p$ is the second-smallest eigenvalue of the
  normalized Laplacian $L_p = I - \frac{1}{d}A_p$, where $A_p$ is the adjacency matrix
  of $\mathcal{G}_p$ (valid since $\mathcal{G}_p$ is $d$-regular when $\Gamma$ acts freely).
\end{definition}

\begin{remark}[Prime lift and CRT splitting]
  \label{rem:prime_lift}
  For applications to the mod-$24$ harmonic seed families, we set $M_p = 24p$ with
  $p > 3$ prime, so $\gcd(24, p) = 1$.  By the Chinese Remainder Theorem,
  \[
    (\mathbb{Z}/M_p\mathbb{Z})^2 \;\cong\; (\mathbb{Z}/24\mathbb{Z})^2 \times (\mathbb{Z}/p\mathbb{Z})^2.
  \]
  The factor $(\mathbb{Z}/24\mathbb{Z})^2$ carries the mod-$24$ resonance structure (fixed
  attractor families, parity layers) and is independent of $p$.
  The factor $(\mathbb{Z}/p\mathbb{Z})^2 \cong \mathbb{F}_p^2$ carries the large-scale
  geometric structure studied in this section.  All spectral and diameter
  statements below refer to the $\mathbb{F}_p^2$ factor.
\end{remark}

\subsection{Harmonic Structure Definitions}
\label{sec:harmonic:foliation}

\begin{definition}[Harmonic foliation]
  \label{def:foliation}
  A \emph{harmonic foliation} of the finite set $V = \mathbb{F}_p^2$ is a
  nontrivial $\Gamma$-invariant partition
  $V = \bigsqcup_i F_i$ (nontrivial means $|V| > |\{F_i\}| > 1$).

  Examples: fixed-parity sublattices; fixed coordinate lines; diagonal cosets;
  fixed mod-$q$ residue classes for $q \mid p-1$.
\end{definition}

\begin{definition}[Harmonic reducibility]
  \label{def:reducibility}
  The pair $(V, \Gamma)$ is \emph{harmonically reducible} if a harmonic foliation
  exists; otherwise it is \emph{harmonically irreducible}.
\end{definition}

\begin{remark}
  Harmonic reducibility of $(V, \Gamma)$ is equivalent to the linear group
  $\rho_p(G) \leq \mathrm{GL}_2(\mathbb{F}_p)$ preserving a nontrivial subspace of
  $\mathbb{F}_p^2$, i.e., the action of $\rho_p(G)$ on $\mathbb{F}_p^2$ is
  \emph{reducible} in the representation-theoretic sense.
\end{remark}

\subsection{Corrected Core Structure Lemma}
\label{sec:harmonic:core}

The following replaces the erroneous non-amenability claim from earlier drafts.
Every finite group is amenable; the correct dichotomy lives at the level of the
\emph{infinite} group $G \leq \mathrm{SL}_2(\mathbb{Z})$.

\begin{lemma}[Tits Alternative, \cite{tits1972free}]
  \label{lem:tits}
  Let $G \leq \mathrm{SL}_2(\mathbb{Z})$ be finitely generated.
  Then exactly one of the following holds:
  \begin{enumerate}
    \item[\emph{(A)}] $G$ is virtually solvable.
    \item[\emph{(B)}] $G$ contains a free subgroup $F_2$ on two generators.
  \end{enumerate}
\end{lemma}

\begin{proof}
  $\mathrm{SL}_2(\mathbb{Z}) \leq \mathrm{GL}_2(\mathbb{Q})$, and $\mathbb{Q}$ has
  characteristic zero, so the Tits Alternative for finitely generated linear
  groups over characteristic-zero fields~\cite{tits1972free} applies directly.
\end{proof}

\begin{remark}[Spectral gap, not amenability]
  \label{rem:amenability}
  Every finite group is amenable; each $\mathrm{SL}_2(\mathbb{F}_p)$ is therefore
  trivially amenable.  Amenability carries no information distinguishing
  expansion behavior among the Cayley graphs $\{\mathcal{G}_p\}$.
  The correct invariant is the \emph{uniform spectral gap} of the family:
  $\inf_p \gamma_p > 0$ (expansion) versus $\gamma_p \to 0$ (non-expansion).
  This is a property of the \emph{family} $\{\mathcal{G}_p\}$, controlled by the
  structure of the infinite lifting group $G \leq \mathrm{SL}_2(\mathbb{Z})$,
  not by amenability of individual finite quotients.
\end{remark}

\subsection{Harmonic Expansion Dichotomy}
\label{sec:harmonic:dichotomy}

\begin{definition}[Zariski-dense subgroup]
  \label{def:zariski}
  $G \leq \mathrm{SL}_2(\mathbb{Z})$ is \emph{Zariski-dense} in $\mathrm{SL}_2$
  if its Zariski closure in $\mathrm{SL}_2(\mathbb{C})$ (in the Zariski topology of
  $\mathbb{A}^4$) equals $\mathrm{SL}_2(\mathbb{C})$.
  Equivalently, $G$ is not contained in any proper algebraic subgroup of
  $\mathrm{SL}_2$ over $\overline{\mathbb{Q}}$.
\end{definition}

\begin{lemma}[Zariski density $\Leftrightarrow$ case (B)]
  \label{lem:zariski_free}
  $G \leq \mathrm{SL}_2(\mathbb{Z})$ is Zariski-dense if and only if $G$ is not
  virtually solvable (equivalently, case~\emph{(B)} of Lemma~\ref{lem:tits} holds).
\end{lemma}

\begin{proof}
  Virtually solvable subgroups of $\mathrm{SL}_2$ are contained in a Borel subgroup
  (triangularizable over $\overline{\mathbb{Q}}$), hence have Zariski closure
  $\subsetneq \mathrm{SL}_2$.  Conversely, if $G$ is not virtually solvable it
  contains $F_2$; by the Tits proof, $F_2 \leq \mathrm{SL}_2(\mathbb{Z})$ is
  Zariski-dense.
\end{proof}

\begin{theorem}[Harmonic Expansion Dichotomy]
  \label{thm:expansion}
  Let $\Gamma \subset \mathrm{SL}_2(\mathbb{Z})$ be a finite symmetric generating set,
  $G = \langle \Gamma \rangle$.  For each prime $p > 3$, let
  $\mathcal{G}_p = \mathrm{Cay}(\rho_p(G), \rho_p(\Gamma))$ with spectral gap $\gamma_p$.
  Then:
  \begin{enumerate}
    \item[\emph{(A)}] \textbf{Virtually solvable case.}
      If $G$ is virtually solvable, then $G$ has polynomial growth, and the family
      $\{\mathcal{G}_p\}$ is \emph{not} an expander: $\gamma_p \to 0$ as $p \to \infty$.

    \item[\emph{(B)}] \textbf{Zariski-dense / free-core case.}
      If $G$ is Zariski-dense in $\mathrm{SL}_2$, then:
      \begin{enumerate}
        \item[\emph{(i)}] $G$ contains a free subgroup $F_2$ (Lemma~\ref{lem:tits}).
        \item[\emph{(ii)}] $\rho_p(G) = \mathrm{SL}_2(\mathbb{F}_p)$ for all but finitely
          many primes $p$ (strong approximation,~\cite{matthews1984congruence}).
        \item[\emph{(iii)}] There exists $\varepsilon > 0$ such that
          $\gamma_p \geq \varepsilon$ for all sufficiently large $p$:
          the family $\{\mathcal{G}_p\}$ is an \emph{expander family}
          (\cite{bourgain2008spectral}).
        \item[\emph{(iv)}] $\mathrm{diam}(\mathcal{G}_p) = O(\log p)$.
      \end{enumerate}
  \end{enumerate}
\end{theorem}

\begin{proof}
  \emph{Case~(A).}
  Virtual solvability implies polynomial growth (Milnor--Wolf theorem).
  Polynomial-growth groups satisfy Kesten's criterion for amenability, from which
  $\gamma_p \to 0$ as the quotient grows.

  \emph{Case~(B)-(i).} Lemma~\ref{lem:tits}.
  The free subgroup $F_2 \leq G$ alone does \emph{not} imply expansion of the
  finite quotients.  Expander families are a property of finite Cayley-graph
  sequences; free subgroup existence in the lifting group $G$ is necessary but
  not sufficient.  Uniform spectral gap of $\{\mathcal{G}_p\}$ requires the
  arithmetic input provided by items~(ii) and~(iii): surjectivity of $\rho_p$
  (strong approximation) and the product-growth machinery of
  Bourgain--Gamburd~\cite{bourgain2008spectral}.

  \emph{Case~(B)-(ii).} Since $G$ is Zariski-dense, the strong-approximation theorem
  of Matthews--\allowbreak{}Vaserstein--\allowbreak{}Weisfeiler~\cite{matthews1984congruence}
  implies $\rho_p(G) = \mathrm{SL}_2(\mathbb{F}_p)$ for all but finitely many $p$.

  \emph{Case~(B)-(iii).}
  We invoke Bourgain--Gamburd~\cite{bourgain2008spectral} under the following
  explicit hypotheses, all of which hold in our setting:
  \begin{itemize}
    \item[(H1)] $\Gamma \subset \mathrm{SL}_2(\mathbb{Z})$ is a \emph{fixed} finite
      symmetric generating set (independent of $p$).
    \item[(H2)] $G = \langle \Gamma \rangle$ is \emph{Zariski-dense} in
      $\mathrm{SL}_2$ over $\overline{\mathbb{Q}}$ (Definition~\ref{def:zariski}).
    \item[(H3)] For all sufficiently large primes $p$,
      $\rho_p(G) = \mathrm{SL}_2(\mathbb{F}_p)$ (established in item~(ii) via
      strong approximation).
  \end{itemize}
  Under (H1)--(H3), Bourgain--Gamburd guarantees the existence of $\varepsilon > 0$
  such that $\gamma_p \geq \varepsilon$ for all sufficiently large $p$:
  the family $\{\mathcal{G}_p\}$ is an expander family.

  \emph{Case~(B)-(iv).} From (iii), the spectral gap $\gamma_p \geq \varepsilon > 0$.
  A $d$-regular expander on $n$ vertices has diameter $O(\log_d n)$
  (standard; see e.g.~\cite{lubotzky1994discrete}).
  Here $n = |\mathrm{SL}_2(\mathbb{F}_p)| = p(p^2 - 1) = O(p^3)$, giving
  $\mathrm{diam}(\mathcal{G}_p) = O(\log p)$.
\end{proof}

\begin{remark}[No entropy; correct observables]
  \label{rem:entropy}
  Earlier drafts claimed $h_M \sim 1/M$ as a ``harmonic entropy.''
  This is incorrect: the topological entropy of any finite deterministic system
  is zero.  The correct observable characterizing expansion is the spectral gap
  $\gamma_p$ (or equivalently the Cheeger constant via the discrete
  Cheeger--Buser inequality), and the correct scaling observable is
  $\mathrm{diam}(\mathcal{G}_p)$.  These replace all entropy claims throughout.
\end{remark}

\subsection{No Continuous-Time Obstruction --- Corrected Statement}
\label{sec:ncto}

\begin{definition}[Polynomial-volume scaling limit]
  \label{def:pvol_limit}
  A family $\{(X_p, d_p)\}_{p \to \infty}$ of metric spaces has a
  \emph{polynomial-volume scaling limit} if there exist scale factors $s_p > 0$
  such that $(X_p, d_p/s_p)$ converges, in the pointed Gromov--Hausdorff sense,
  to a metric measure space $(Y, d_Y, \mu)$ whose volume satisfies
  \[
    \mu\bigl(B(y, r)\bigr) \;\leq\; C r^n
  \]
  for some constants $C, n > 0$, all $y \in Y$, and all $r > 0$.

  This class includes: connected Riemannian manifolds (Bishop--Gromov), flat tori,
  and Carnot groups (Pansu's volume formula).  It is the natural target for
  rescaling limits of polynomial-growth groups.  Fractal and $p$-adic spaces are
  excluded.

  For our setting: $X_p$ is the vertex set of $\mathcal{G}_p$,
  $d_p = T_\Gamma$ is the graph distance (reachability time),
  and we ask whether any sequence $s_p$ produces such a limit.
\end{definition}

\begin{lemma}[Growth obstruction]
  \label{lem:growth_obs}
  Let $\{(X_p, d_p)\}_{p \geq p_0}$ be a family of connected $d$-regular graphs
  with uniform spectral gap $\gamma_p \geq \varepsilon > 0$.  Then:
  \begin{enumerate}
    \item[\emph{(i)}] \textbf{Expander ball growth.} There exists
      $c = c(\varepsilon, d) > 0$ such that
      \[
        |B(x, k)| \;\geq\; (1 + c)^{\,k}
      \]
      for all vertices $x \in X_p$ and all $1 \leq k \leq \tfrac{1}{2}\mathrm{diam}(X_p)$.
    \item[\emph{(ii)}] \textbf{No polynomial-volume GH limit.} For every sequence
      $s_p > 0$, the rescaled family $(X_p,\, d_p / s_p)$ does \emph{not}
      Gromov--Hausdorff converge to any metric measure space with polynomial
      volume growth.
  \end{enumerate}
\end{lemma}

\begin{proof}
  \emph{(i).}  Uniform spectral gap $\gamma_p \geq \varepsilon$ implies, via the
  discrete Cheeger--Buser inequality, that the edge-expansion constant
  $h(X_p) \geq \varepsilon/2$.  Standard ball-growth estimates for graphs with
  positive Cheeger constant give exponential growth with $c \geq \varepsilon/(2d)$.

  \emph{(ii).}  Suppose $(X_p, d_p/s_p) \xrightarrow{GH} (Y, d_Y, \mu)$ where
  $\mu(B(y, r)) \leq C r^n$ for some $C, n > 0$.  For GH convergence to be
  nondegenerate we need $s_p = \Theta(\mathrm{diam}(X_p))$.  For any $r > 0$,
  the vertex count within rescaled distance $r$ satisfies
  $|B(x, r s_p)| \geq (1+c)^{r s_p}$, which grows without bound in $p$,
  contradicting polynomial growth of $\mu$.
\end{proof}

\begin{theorem}[Generator Growth Dichotomy]
  \label{thm:growth_dichotomy}
  Let $\Gamma \subset \mathrm{SL}_2(\mathbb{Z})$ be finite symmetric,
  $G = \langle\Gamma\rangle$.  Exactly one of the following holds:
  \begin{enumerate}
    \item[\emph{(A)}] \textbf{Polynomial growth regime.}
      $G$ is virtually solvable.  Ball volumes in $\mathcal{G}_p$ grow polynomially
      in radius, and the family admits a polynomial-volume scaling limit.

    \item[\emph{(B)}] \textbf{Exponential growth regime.}
      $G$ is Zariski-dense.  There exists $c = c(\Gamma) > 0$ such that for all
      sufficiently large primes $p$,
      \[
        \bigl|B_{\mathcal{G}_p}(x, r)\bigr| \;\geq\; e^{\,cr}
      \]
      for all vertices $x$ and all $1 \leq r \leq C\log p$.
      The family admits \emph{no} polynomial-volume scaling limit.
  \end{enumerate}
  The growth regime is determined solely by the Zariski closure of
  $G$ in $\mathrm{SL}_2$.
\end{theorem}

\begin{proof}
  \emph{Case~(A).}
  Virtual solvability in $\mathrm{SL}_2(\mathbb{Z})$ implies virtual abelianness:
  the proper algebraic subgroups of $\mathrm{SL}_2$ are Borel (unipotent) or
  split torus, both virtually abelian.  Virtually abelian groups have polynomial
  growth (Milnor--Wolf theorem).  The congruence quotients $G_p = \rho_p(G)$
  are abelian of bounded rank; their Cayley graphs have polynomial ball growth
  (radius-$r$ ball has size $O(r^{\mathrm{rank}})$), and rescale by
  $s_p = |G_p|^{1/\mathrm{rank}}$ to flat Riemannian tori, which have
  polynomial volume growth.  A polynomial-volume scaling limit exists.

  \emph{Case~(B).}
  $G$ Zariski-dense $\Rightarrow$ by Theorem~\ref{thm:expansion}(B)(iii),
  $\gamma_p \geq \varepsilon > 0$ for all large $p$.
  The discrete Cheeger--Buser inequality gives $h(\mathcal{G}_p) \geq \varepsilon/2$,
  whence $|B(x, r)| \geq (1 + c')^r$ with $c' = \varepsilon/(2d)$; set
  $c = \log(1 + c') > 0$.  This is the exponential lower bound.

  For the no-limit claim: suppose $(Y, d_Y, \mu)$ with
  $\mu(B(y,r)) \leq Cr^n$ is a GH limit of the rescaled family.
  By Theorem~\ref{thm:expansion}(B)(iv), $s_p = \Theta(\log p)$.
  For any fixed $r > 0$,
  \[|B(x, r s_p)| \;\geq\; e^{cr s_p} \;\to\; \infty \quad (p \to \infty),\]
  contradicting polynomial volume growth of~$\mu$.
\end{proof}

\begin{corollary}[No Continuous-Time Obstruction]
  \label{cor:ncto}
  In the exponential growth regime (Case~\emph{(B)} of
  Theorem~\ref{thm:growth_dichotomy}), the family
  $\{(\mathcal{G}_p, T_\Gamma)\}$ admits no polynomial-volume scaling limit.
  In particular, it admits no Riemannian and no sub-Riemannian (Carnot) scaling
  limit.
\end{corollary}

\begin{remark}[Role of spectral gap]
  \label{rem:gap_role}
  The spectral gap $\gamma_p \geq \varepsilon > 0$ (Bourgain--Gamburd) is used in
  Case~(B) as an intermediate step: Zariski-dense $\Rightarrow$ uniform gap
  $\Rightarrow$ exponential ball growth.  It is not the structural invariant
  of the theorem.  The invariant is the \emph{growth class} of the generator
  group, which cleanly separates polynomial-volume limits from non-existence of
  any such limit, across both Riemannian and sub-Riemannian target spaces.
\end{remark}

\subsection{NCTO as a Generator-Relative Obstruction Principle}
\label{sec:ncto:gen}

Theorem~\ref{thm:growth_dichotomy} characterises growth class from $\Gamma$ alone.
The implication chain below extracts the spectral and geometric consequences
in computable form.

\begin{definition}[Generator obstruction tag]
  \label{def:ncto_tag}
  A finite symmetric $\Gamma \subset \mathrm{SL}_2(\mathbb{Z})$ carries the
  \emph{NCTO tag}
  \[
    \tau(\Gamma) \;=\;
    \begin{cases}
      \texttt{ncto-blocked} & \text{if } \langle\Gamma\rangle
                              \text{ is Zariski-dense in } \mathrm{SL}_2, \\
      \texttt{ok}           & \text{if } \langle\Gamma\rangle
                              \text{ is virtually solvable.}
    \end{cases}
  \]
  By Lemma~\ref{lem:tits}, exactly one case applies to any finitely generated
  $G \leq \mathrm{SL}_2(\mathbb{Z})$.
\end{definition}

\begin{theorem}[NCTO, Generator-Relative Implications]
  \label{thm:ncto_equiv}
  Let $\Gamma \subset \mathrm{SL}_2(\mathbb{Z})$ be finite symmetric,
  $G = \langle\Gamma\rangle$.  Then:
  \[
    \text{(i)} \;\Rightarrow\; \text{(ii)} \;\Rightarrow\; \text{(iii)}
    \;\Rightarrow\; \text{(iv)},
  \]
  where:
  \begin{enumerate}
    \item[\emph{(i)}]   $\tau(\Gamma) = \texttt{ncto-blocked}$ ($G$ Zariski-dense).
    \item[\emph{(ii)}]  There exists $\varepsilon = \varepsilon(\Gamma) > 0$ such that
                        $\gamma_p \geq \varepsilon$ for all sufficiently large $p$.
    \item[\emph{(iii)}] $\mathrm{diam}(\mathcal{G}_p) = O(\log p)$.
    \item[\emph{(iv)}]  $\{(\mathcal{G}_p, T_\Gamma)\}$ has no polynomial-volume
                        scaling limit (Definition~\ref{def:pvol_limit}).
  \end{enumerate}
  Moreover, $\varepsilon(\Gamma)$ depends only on the Zariski closure of $G$;
  for $\Gamma = \{\sigma^{\pm 1}, w^{\pm 1}\}$ it is made explicit by
  Bourgain--Gamburd~\cite{bourgain2008spectral}.
\end{theorem}

\begin{proof}
  \emph{(i) $\Rightarrow$ (ii).}
  Theorem~\ref{thm:expansion}(B)(iii): Bourgain--Gamburd under (H1)--(H3).

  \emph{(ii) $\Rightarrow$ (iii).}
  A $d$-regular expander on $n$ vertices has diameter $O(\log n)$
  (see e.g.~\cite{lubotzky1994discrete}).
  Here $n = |\mathrm{SL}_2(\mathbb{F}_p)| = O(p^3)$, giving $O(\log p)$.

  \emph{(iii) $\Rightarrow$ (iv).}
  Logarithmic diameter implies exponential ball growth (Lemma~\ref{lem:growth_obs}(i)),
  which prevents GH convergence to any polynomial-volume space
  (Lemma~\ref{lem:growth_obs}(ii)).
\end{proof}

\begin{remark}[Converses fail]
  \label{rem:ncto_conv}
  The implications (ii)$\Rightarrow$(iii) and (iii)$\Rightarrow$(iv) are strict.
  (ii) $\not\Leftarrow$ (iii): there exist bounded-degree graph families with
  diameter $O(\log n)$ and spectral gap $\gamma_p \to 0$; logarithmic diameter
  alone does not force a spectral gap lower bound.
  (i) $\not\Leftarrow$ (iv): Definition~\ref{def:pvol_limit} already includes
  Carnot (sub-Riemannian) limits, so virtually nilpotent groups (whose Pansu
  limits are Carnot groups) are correctly placed on the ``limit exists'' side.
  However, (iv) alone does not force Zariski-density without the growth
  classification of Theorem~\ref{thm:growth_dichotomy}: one must also rule out
  other reasons for non-convergence (e.g.\ non-compactness of the limit).
\end{remark}

\begin{remark}[QA obstruction types]
  \label{rem:ncto_type}
  In the QA generator-certificate framework, the NCTO tag classifies generators by
  the \emph{type} of geometric obstruction they carry:
  \begin{itemize}
    \item \textbf{Type~I (topological):} connectivity or reachability obstruction;
      the Cayley graph is disconnected or has isolated vertices.
    \item \textbf{Type~II (metric / PDI):} route-reducing obstruction; the generator
      is tagged \texttt{parity-block}, reducing multi-path states and lowering the
      Path Diversity Index.
    \item \textbf{Type~III (spectral / NCTO):} expansion obstruction.
      Tag: $\tau(\Gamma) = \texttt{ncto-blocked}$.
      Effect: uniform gap $\varepsilon(\Gamma) > 0$, no GH Riemannian limit.
  \end{itemize}
  Unlike Type~I and Type~II obstructions, which are local (they depend on individual
  generator edges), the Type~III obstruction is \emph{global and family-level}: it
  is a property of the entire congruence quotient family $\{\mathcal{G}_p\}$,
  determined by the Zariski closure of $G = \langle\Gamma\rangle$ in $\mathrm{SL}_2$.

  The tag $\tau(\Gamma)$ is finitely \emph{decidable} from $\Gamma$: algorithms
  for testing virtual solvability in $\mathrm{SL}_2(\mathbb{Z})$ are known.
  In practice, reduction modulo several primes provides strong numerical evidence
  (if $\rho_p(\langle\Gamma\rangle) = \mathrm{SL}_2(\mathbb{F}_p)$ for multiple
  primes, Zariski-density is highly indicated), but a single prime test does not
  constitute a proof of the Bourgain--Gamburd hypotheses across all large~$p$.
\end{remark}

\subsection{Generator Specification and Correction}
\label{sec:harmonic:generators}

\begin{definition}[Prime-lifted tonal action]
  \label{def:tonal_action}
  Let $N = 24$ be the chromatic period of the mod-$24$ QA harmonic seed system,
  and let $p > 3$ be prime with $M_p = 24p$.  The \emph{prime-lifted tonal action}
  is the action of $G_0 = \langle \sigma, w \rangle = \mathrm{SL}_2(\mathbb{Z})$ on
  $(\mathbb{Z}/M_p\mathbb{Z})^2$ via entry-wise reduction mod $M_p$.

  The two basic tonal operations are identified as follows:
  \begin{itemize}
    \item \textbf{Semitone transposition}: $\sigma = \bigl[\begin{smallmatrix}1&1\\0&1\end{smallmatrix}\bigr]$
      shifts the ``base'' coordinate by one step mod $M_p$, corresponding to ascending
      one semitone in the $24$-tone chromatic lattice.
    \item \textbf{Melodic inversion}: $w = \bigl[\begin{smallmatrix}0&1\\-1&0\end{smallmatrix}\bigr]$
      (Weyl element) swaps and negates coordinates, corresponding to inversion of a
      pitch-class set around the central axis.
  \end{itemize}

  By the CRT splitting $(\mathbb{Z}/M_p\mathbb{Z})^2 \cong
  (\mathbb{Z}/24\mathbb{Z})^2 \times \mathbb{F}_p^2$ (Remark~\ref{rem:prime_lift}),
  the mod-$24$ factor records the harmonic seed structure (attractor families, parity
  layers), while the $\mathbb{F}_p^2$ factor is the Cayley-graph domain analyzed by
  Bourgain--Gamburd.  As $p \to \infty$, the lattice resolution increases, and the
  spectral geometry of $\mathcal{G}_p$ (Theorem~\ref{thm:expansion}) governs the
  large-scale tonal structure.
\end{definition}

Earlier drafts used generators $\sigma, \mu, \lambda$ where
$\mu = \bigl[\begin{smallmatrix}0&1\\1&0\end{smallmatrix}\bigr]$ has
$\det \mu = -1 \notin \mathrm{SL}_2(\mathbb{Z})$.
The corrected generator set, all lying in $\mathrm{SL}_2(\mathbb{Z})$, is:
\[
  \sigma \;=\; \begin{pmatrix} 1 & 1 \\ 0 & 1 \end{pmatrix}, \qquad
  w \;=\; \begin{pmatrix} 0 & 1 \\ -1 & 0 \end{pmatrix}, \qquad
  \lambda_k \;=\; \begin{pmatrix} k & 0 \\ 0 & k^{-1} \end{pmatrix}
\]
where $k \in \mathbb{Z}$ with $k \neq \pm 1$ (so $\lambda_k$ is not unipotent) and
$\gcd(k, M_p) = 1$ (so $\lambda_k$ is invertible mod $M_p$).

Note: $w$ is the standard Weyl element of order~4 in $\mathrm{SL}_2(\mathbb{Z})$;
it replaces $\mu$ and satisfies $\det w = 1$.  In the classical presentation of the
modular group, set
\[
  T \;:=\; \sigma \;=\; \begin{pmatrix}1&1\\0&1\end{pmatrix}, \qquad
  S \;:=\; w^{-1} \;=\; \begin{pmatrix}0&-1\\1&0\end{pmatrix};
\]
then $\mathrm{SL}_2(\mathbb{Z}) = \langle S, T \mid S^4 = I,\ (ST)^3 = S^2 \rangle$
(so $\langle \sigma, w \rangle = \mathrm{SL}_2(\mathbb{Z})$).
This identification makes the Zariski-density of $\langle \sigma, w \rangle$
immediate: $\mathrm{SL}_2(\mathbb{Z})$ is a lattice in $\mathrm{SL}_2(\mathbb{R})$
and hence Zariski-dense.
Adding $\lambda_k$ with $k \neq \pm 1$ preserves Zariski-density (the group only grows)
while ensuring the generator set is not contained in a unipotent or torus subgroup.

\begin{lemma}[Zariski density of $\langle \sigma, w \rangle$]
  \label{lem:generators_dense}
  The group $G_0 = \langle \sigma, w \rangle = \mathrm{SL}_2(\mathbb{Z})$
  is Zariski-dense in $\mathrm{SL}_2$.
\end{lemma}

\begin{proof}
  $\mathrm{SL}_2(\mathbb{Z})$ is a lattice in $\mathrm{SL}_2(\mathbb{R})$, hence
  Zariski-dense in $\mathrm{SL}_2(\mathbb{C})$.  Setting $T = \sigma$ and
  $S = w^{-1}$, the standard presentation
  $\langle S, T \mid S^4 = I,\ (ST)^3 = S^2 \rangle$
  is verified in \S\ref{sec:harmonic:generators}.
\end{proof}

\begin{remark}
  Adding $\lambda_k = \bigl[\begin{smallmatrix}k&0\\0&k^{-1}\end{smallmatrix}\bigr]$
  for $k \in \mathbb{Z} \setminus \{0, \pm 1\}$ with $\gcd(k,6)=1$ preserves
  Zariski-density (the group only grows).  It additionally ensures the generator set
  is not confined to a unipotent or torus subgroup, which is relevant when $\lambda_k$
  is used in the prime-lifted tonal action
  (Definition~\ref{def:tonal_action}).
\end{remark}

\subsection{Reflective versus Rotational Tonal Inversion}
\label{sec:harmonic:inversion_type}

We distinguish two inversion types that arise in tonal systems.
Call $r: p \mapsto -p$ (pitch-class reversal on $\mathbb{Z}/n\mathbb{Z}$) a
\emph{reflective} inversion, and call the Weyl rotation
$(b,e) \mapsto (e,-b)$ on $\mathbb{Z}^2$ a \emph{rotational} inversion.
Both are called ``inversion'' in common usage, but they act on spaces of
different dimension and land on opposite sides of
Theorem~\ref{thm:growth_dichotomy}.

\begin{definition}[Reflective inversion]
  \label{def:refl_inv}
  A \emph{reflective inversion} on a cyclic pitch space $\mathbb{Z}/n\mathbb{Z}$
  is an automorphism $r$ satisfying
  \[
    r^2 = \mathrm{id}, \qquad r\sigma r^{-1} = \sigma^{-1}.
  \]
  The standard example is $r(p) = -p \pmod{n}$ (pitch-class inversion).
  The symmetry group $\langle \sigma, r \rangle$ is the dihedral group $D_n$.
\end{definition}

\begin{definition}[Rotational inversion]
  \label{def:rot_inv}
  A \emph{rotational inversion} compatible with $\sigma$ is an automorphism
  $w \in \mathrm{SL}_2(\mathbb{Z})$ with $w^2 = -\mathrm{id}$ and $\det(w) = 1$.
  Setting $S := w^{-1}$ and $T := \sigma$, we require the modular relation
  \[
    S^4 = I, \qquad (ST)^3 = S^2.
  \]
  These are the standard $\mathrm{SL}_2(\mathbb{Z})$ presentation relations
  (Lemma~\ref{lem:harm_gen}); they tie $w$ to the translation direction $\sigma$.

  \emph{Remark.}  For the canonical $\sigma =
  \bigl[\begin{smallmatrix}1&1\\0&1\end{smallmatrix}\bigr]$ and
  $w = \bigl[\begin{smallmatrix}0&1\\-1&0\end{smallmatrix}\bigr]$,
  \[
    w\sigma w^{-1} = \begin{pmatrix}1&0\\-1&1\end{pmatrix} =: \sigma_{\downarrow},
  \]
  so $w$ conjugates $\sigma$ (upper unipotent) to $\sigma_{\downarrow}$ (lower
  unipotent): a harmonic role exchange, not inversion to $\sigma^{-1}$.
  The modular relation $(ST)^3 = S^2$ is the correct algebraic tie.
  The Weyl element $w = \bigl[\begin{smallmatrix}0&1\\-1&0\end{smallmatrix}\bigr]$
  (acting as $(b,e) \mapsto (e,-b)$) is the canonical example.
\end{definition}

\begin{lemma}[Harmonic lattice generation]
  \label{lem:harm_gen}
  Let $\sigma = \bigl[\begin{smallmatrix}1&1\\0&1\end{smallmatrix}\bigr]$ and
  $w = \bigl[\begin{smallmatrix}0&1\\-1&0\end{smallmatrix}\bigr]$.
  Then $w^2 = -\mathrm{id}$, $\det(w) = 1$, and
  $\langle \sigma, w \rangle = \mathrm{SL}_2(\mathbb{Z})$.
\end{lemma}

\begin{proof}
  Set $S := w^{-1} = \bigl[\begin{smallmatrix}0&-1\\1&0\end{smallmatrix}\bigr]$
  and $T := \sigma$.
  Direct matrix computation gives $S^4 = I$ and $(ST)^3 = S^2 = -I$
  (see \S\ref{sec:harmonic:generators}).
  The presentation $\langle S, T \mid S^4 = I,\ (ST)^3 = S^2 \rangle$
  is a classical generator theorem for $\mathrm{SL}_2(\mathbb{Z})$
  \cite{lubotzky1994discrete}.
  Hence $\langle \sigma, w \rangle = \langle T, S \rangle = \mathrm{SL}_2(\mathbb{Z})$.
\end{proof}

\begin{proposition}[Tonal Inversion Dichotomy]
  \label{prop:tonal_dict}
  Let $\mathcal{T}$ be a tonal system closed under semitone transposition $\sigma$.
  \begin{enumerate}
    \item[\emph{(A)}] If $\mathcal{T}$ admits a \emph{reflective} inversion $r$
      on a $1$-dimensional pitch space $\mathbb{Z}/n\mathbb{Z}$, then
      $G_\mathcal{T} = \langle\sigma, r\rangle \cong D_n$ is virtually abelian:
      polynomial ball growth, Case~(A) of Theorem~\ref{thm:growth_dichotomy}.
    \item[\emph{(B)}] If $\mathcal{T}$ admits a $\sigma$-compatible rotational
      inversion $w$ (Definition~\ref{def:rot_inv}), then by
      Lemma~\ref{lem:harm_gen},
      $\langle\sigma,w\rangle \cong \mathrm{SL}_2(\mathbb{Z})$;
      in particular $G_\mathcal{T}$ is Zariski-dense:
      exponential ball growth, Case~(B) of
      Theorem~\ref{thm:growth_dichotomy}.
  \end{enumerate}
  The regime is determined by the inversion type: reflective ($r^2=\mathrm{id}$,
  $1$-dimensional) yields Case~(A); $\sigma$-compatible rotational
  ($S^4=I$, $(ST)^3=S^2$) yields Case~(B).
\end{proposition}

\begin{proof}
  \emph{Case~(A).}  $D_n$ is virtually abelian;
  Theorem~\ref{thm:growth_dichotomy}(A) applies.

  \emph{Case~(B).}  By Lemma~\ref{lem:harm_gen},
  $\langle\sigma, w\rangle = \mathrm{SL}_2(\mathbb{Z})$.
  Apply Theorem~\ref{thm:growth_dichotomy}(B) and Corollary~\ref{cor:ncto}.
\end{proof}

\begin{remark}[Dimensional elevation]
  \label{rem:dim_elevation}
  Reflective inversion is an order-$2$ isometry of the pitch circle $S^1$; it
  preserves the one-dimensional geometry.  Rotational inversion cannot be realised
  as an isometry of $S^1$: it requires a $2$-dimensional harmonic lattice.
  The passage from reflective to rotational inversion is therefore a dimensional
  elevation from $S^1$ to the hyperbolic plane $\mathbb{H}$ (on which
  $\mathrm{SL}_2(\mathbb{Z})$ acts via M\"{o}bius transformations), and this
  elevation is precisely what forces the exponential growth regime.
\end{remark}

\subsection{Worked Examples --- Corrected}
\label{sec:harmonic:examples}

\paragraph{Example A --- Virtually Abelian (Case A).}
Let $\Gamma = \{\pm\sigma\}$ where $\sigma = \bigl[\begin{smallmatrix}1&1\\0&1\end{smallmatrix}\bigr]$.
Then $G = \langle \sigma \rangle \cong \mathbb{Z}$ is abelian, hence virtually solvable.

The Cayley graph $\mathcal{G}_p = \mathrm{Cay}(\mathbb{Z}/p\mathbb{Z}, \{\pm 1\})$ is a
cycle of length $p$.  Its spectral gap $\gamma_p = 1 - \cos(2\pi/p) \sim 2\pi^2/p^2 \to 0$.
Its diameter is $\lfloor p/2 \rfloor \sim p/2$.

This family is \emph{not} an expander: $\gamma_p \to 0$.  Rescaling by $s_p = p$
gives convergence to the circle $S^1$: a smooth continuous limit exists.
\emph{Continuous-time behavior is possible.}

\paragraph{Example B (Case B: Zariski-dense).}
Let $\Gamma = \{\sigma^{\pm 1}, w^{\pm 1}\}$ as in Definition~\ref{def:tonal_action}.
Then $G = \mathrm{SL}_2(\mathbb{Z})$ (Zariski-dense lattice).

By Bourgain--Gamburd~\cite{bourgain2008spectral} (Theorem~\ref{thm:expansion}(B)(iii)),
the Cayley graphs $\mathcal{G}_p = \mathrm{Cay}(\mathrm{SL}_2(\mathbb{F}_p), \rho_p(\Gamma))$
form an expander family with uniform gap $\gamma_p \geq \varepsilon > 0$.

Diameter: $\mathrm{diam}(\mathcal{G}_p) = O(\log p)$.
Ball volume: exponential in graph distance ($|B(x,r)| \geq e^{cr}$,
Theorem~\ref{thm:growth_dichotomy}(B)).
No rescaling produces any polynomial-volume limit (Corollary~\ref{cor:ncto}).
\emph{Harmonic motion is geometrically rigid.}

\subsection{Structural Summary}
\label{sec:harmonic:summary}

\[
\renewcommand{\arraystretch}{1.3}
\begin{array}{ccccc}
  G \text{ virt.\ solvable} & \Longleftrightarrow &
  \text{poly.\ ball growth} & \Rightarrow &
  \text{poly-vol limit exists} \\
  G \text{ Zariski-dense} & \Longleftrightarrow &
  \text{exp.\ ball growth} & \Rightarrow &
  \text{no poly-vol limit}
\end{array}
\]

The growth class (polynomial vs.\ exponential) is the structural invariant,
determined entirely by the Zariski closure of $G \leq \mathrm{SL}_2(\mathbb{Z})$.
Spectral gap ($\gamma_p$) is a consequence of Zariski-density, not the primary
object; it is used as an intermediate tool (Bourgain--Gamburd $\Rightarrow$
Cheeger $\Rightarrow$ exponential growth) rather than the dichotomy's engine.

\subsection{Tonal Consequence}
\label{sec:harmonic:tonal}

The following corollary instantiates the abstract dichotomy in concrete musical terms,
using the generators identified in Definition~\ref{def:tonal_action}.

\begin{corollary}[Tonal Regime Dichotomy]
  \label{cor:tonal}
  Under the prime-lifted tonal action (Definition~\ref{def:tonal_action}):
  \begin{enumerate}
    \item[\emph{(A)}] \textbf{Transposition-only regime.}
      Let $\Gamma = \{\sigma^{\pm 1}\}$ (pure semitone transposition), so
      $G \cong \mathbb{Z}$.
      Restricting to any $\langle\sigma\rangle$-orbit, the Cayley graph is a cycle
      $C_{M_p}$ of length $M_p = 24p$.  Hence $\gamma_p \to 0$ and the rescaled cycle
      converges to $S^1$ (cf.\ Example~A).
      Harmonic motion in this regime is \emph{geometrically smooth}: it embeds in a
      continuous Riemannian flow.

    \item[\emph{(B)}] \textbf{Harmonic inversion regime.}
      Let $\Gamma \supseteq \{\sigma^{\pm 1}, w^{\pm 1}\}$ with $w$ a rotational
      inversion (Definition~\ref{def:rot_inv}).
      Proposition~\ref{prop:tonal_dict}(B) gives $G = \mathrm{SL}_2(\mathbb{Z})$.
      Ball growth in $\{\mathcal{G}_p\}$ is exponential; by
      Corollary~\ref{cor:ncto}, no polynomial-volume scaling limit exists.
      Harmonic motion is \emph{geometrically rigid}: no rescaling of
      reachability time produces a Riemannian polynomial-volume scaling limit.
  \end{enumerate}
\end{corollary}

\begin{proof}
  Case~(A): $G \cong \mathbb{Z}$, virtually abelian;
  Theorem~\ref{thm:growth_dichotomy}(A) and Example~A.
  Case~(B): $G = \mathrm{SL}_2(\mathbb{Z})$ by
  Proposition~\ref{prop:tonal_dict}(B);
  Theorem~\ref{thm:growth_dichotomy}(B) and Corollary~\ref{cor:ncto}.
\end{proof}

\begin{remark}[Musical interpretation]
  The dichotomy is a property of the \emph{generator regime}, not of individual chord
  moves.  A harmonic system whose generator regime is transposition-only has abelian lift
  ($G \cong \mathbb{Z}$); its large-scale geometry is a circle, admitting smooth
  voice-leading flows.
  A system that combines transposition with inversion generates $\mathrm{SL}_2(\mathbb{Z})$,
  whose congruence Cayley graphs have \emph{exponential ball growth}.
  The Generator Growth Dichotomy says: \emph{exponential growth is incompatible
  with any polynomial-volume scaling limit.}
  No rescaling of reachability time can produce a smooth, Riemannian, or
  sub-Riemannian description of the large-scale harmonic motion.
\end{remark}

% ─── Bibliography ─────────────────────────────────────────────────────────────
% All four required entries live in refs_music.bib (same directory):
%   tits1972free, matthews1984congruence, bourgain2008spectral, lubotzky1994discrete
%
% In the master document preamble use:
%   \bibliography{refs,refs_music}
